\documentclass[10pt,a4paper]{amsart}
\usepackage[margin=19mm]{geometry}
\usepackage{amsmath,amssymb,mathtools}
\usepackage[scr=rsfs,cal=euler]{mathalfa}
\usepackage{xfrac}
\usepackage[unicode,hidelinks,bookmarks=false]{hyperref}
\newcommand{\ie}{i.e.\ }
\newcommand{\cf}{cf.\ }
\newcommand{\resp}{respectively\ }
\newcommand{\Subsection}{\subsection}
\providecommand{\nobreakdash}{\nobreak\hbox{-}}
\setlength{\emergencystretch}{2em}
\multlinegap0pt
\usepackage{bm}
\usepackage{bbm}
\usepackage{dsfont}
\usepackage{xspace}
\usepackage{cancel}
\usepackage{mathtools}
\usepackage{amsthm}
\makeatletter
\@ifundefined{lemma}{\newtheorem{lemma}{Lemma}}{}
\@ifundefined{theorem}{\newtheorem{theorem}{Theorem}}{}
\@ifundefined{corollary}{\newtheorem{corollary}[theorem]{Corollary}}{}
\makeatother
\makeatletter
\newcommand{\Ab}{\mathbf{A}}
\newcommand{\Rcal}{\mathcal{R}}
\newcommand{\Tcal}{\mathcal{T}}
\newcommand{\abr}[1]{\left|#1\right|}
\expandafter\ifx\csname proof\endcsname\relax
\newenvironment{proof}{\par\noindent{\bf Proof\ }}{\hfill\BlackBox\par\medskip}
\fi
\newcommand{\rbr}[1]{\left(#1\right)}
\newcommand{\sbr}[1]{\left[#1\right]}
\makeatother
\title[Stochastic Rounding Increases Small Singular Values]{Stochastic Rounding Increases Small Singular Values}
\author{Linkai Ma, Tingzhou Yu, Petros Drineas}
\date{Source: arXiv:2606.00312v1 (CC BY 4.0)}
\begin{document}
\maketitle

\noindent\emph{Source and licence.} Stochastic Rounding Increases Small Singular Values. Version arXiv:2606.00312v1 (CC BY 4.0), distributed under the Creative Commons Attribution 4.0 International licence, as recorded in the arXiv licence metadata for this exact version and shown on the abstract page https://arxiv.org/abs/2606.00312v1. Full rights notice: licenses/arxiv-2606.00312-CC-BY-4.0.txt.\par\smallskip

\noindent\textit{Editorial scope.} Counted unit: the Corollary whose source label is \texttt{cor:smallcluster\_alter} in the article (source0.tex lines 2796--2805, source label \texttt{cor:smallcluster\_alter}), delivered together with the article's own complete original proof of it (source0.tex lines 2807--2853, one \texttt{proof} environment of 47 lines). No mathematics is written, repaired or re-derived: statement and proof are the article's own text line for line. Required context retained uncounted: eq:alpha-beta (lines 1600--1604); cor:smallcluster (lines 1642--1647); eq:alpha\_tilde (lines 2769--2771); lem:remainder\_alter (lines 2774--2779); thm:smallcluster-main\_alter (lines 2785--2791). Every definition, display and label of those lines is retained unchanged. Omitted from this package and not redistributed: 1--1599, 1605--1641, 1648--2768, 2772--2773, 2780--2784, 2792--2795, 2806--2806, 2854--2860. Nothing inside a retained excerpt is omitted or altered: every retained body is delivered exactly as the article's lines are written, so no omission receipt of any kind accompanies this package. Citations: the retained excerpts carry no citation command at all, so no bibliography entry of the article is transcribed and no reference is redistributed. Reference adaptation: none was needed and none was made; every reference inside the delivered unit resolves inside it, so no unresolved reference or citation is produced. Printed slips kept literal: no printed word, symbol, hypothesis or number of the counted statement or of its proof was altered. Layout and class: the article's own class is \texttt{article} with packages including main, inputenc, fontenc, hyperref, url, booktabs, amsfonts, nicefrac, microtype, xcolor, tikz, bbm, latexsym, amsmath; the wrapper is the lane's pinned \texttt{amsart} editorial wrapper at 10pt on A4, which restates the theorem-like environments the retained text needs and adds the article's own macro definitions that the retained text needs (\texttt{\textbackslash Ab}, \texttt{\textbackslash Rcal}, \texttt{\textbackslash Tcal}, \texttt{\textbackslash abr}, \texttt{\textbackslash rbr}, \texttt{\textbackslash sbr}), with extra wrapper packages (bm, bbm, dsfont, xspace, cancel, mathtools). The delivered numbering is the unit's own. Assets: no retained span contains a figure environment, a graphics-inclusion command, a PDF-page inclusion, a table or a TikZ picture; no image file is copied into this package, the package declares no asset dependency and no image file is redistributed with it. Licence: version arXiv:2606.00312v1 of this article is distributed under the Creative Commons Attribution 4.0 International licence, as recorded in the arXiv licence metadata for this exact version and shown on the abstract page https://arxiv.org/abs/2606.00312v1. A full rights notice is delivered at licenses/arxiv-2606.00312-CC-BY-4.0.txt. Characters: every retained excerpt is pure ASCII, so the delivered engine, XeLaTeX with its default Latin Modern text font, reports no missing character.
\begin{equation}\label{eq:alpha-beta}
    \alpha(\rho) \;:=\; \frac{2}{g}\,C\rbr{\sigma_1(\Ab)\,\rho\sqrt{d} + \rho^2 n},
    \qquad
    \beta(\rho) \;:=\; \frac{2}{g}\rbr{2\,\sigma_1(\Ab)\,\rho\sqrt{nd} + \rho^2 nd}.
\end{equation}

\begin{corollary}\label{cor:smallcluster}
There exist absolute constants $c_1, c_2, c_3, c_g > 0$ such that, if $\rho\sqrt{nd} \le c_1\sigma_k$, $\sigma_1 \le c_2\sigma_k$, $d \ge c_3\sqrt n, g= \lambda_k - \lambda_{k+1} > c_g \lambda_k$, then there is a constant $C > 0$ such that
\begin{equation*}
    \Tcal \ge \sum_{j=k+1}^{d}\nu_j^\downarrow - C\rbr{k(d-k)\,\rho^2 + \frac{k\,n^2 \rho^4}{\sigma_k^2} + \frac{d^{5/2}\rho^3}{\sigma_k}}.
\end{equation*}
\end{corollary}

\begin{equation}\label{eq:alpha_tilde}
    \tilde{\alpha}(\rho):= \; \frac{2}{g}\,C\rbr{\sigma_1(\Ab)\,\rho\sqrt{n} + \rho^2 n}.
\end{equation}

\begin{lemma}[Higher-order remainder alternative]\label{lem:remainder_alter}
Assume $\tilde{\alpha}(\rho) < 1$ and $\beta(\rho) < 1$ with $\tilde{\alpha},\beta$ as in~\eqref{eq:alpha-beta},~\eqref{eq:alpha_tilde}. Then there is an absolute constant $C > 0$ such that
\begin{equation*}
    \abr{\sum_{\ell=3}^{\infty} I_\ell} \le \frac{Cd}{g}\rbr{2(\lambda_1 - \lambda_k) + \pi g}\rbr{\lambda_1 + \frac{g}{2}}\rbr{\frac{\tilde{\alpha}(\rho)^3}{1 - \tilde{\alpha}(\rho)} + 2 e^{-(n+d)}\,\frac{\beta(\rho)^3}{1 - \beta(\rho)}}.
\end{equation*}
\end{lemma}

\begin{theorem}\label{thm:smallcluster-main_alter}
Assume $\tilde{\alpha}(\rho) < 1$ and $\beta(\rho) < 1$ with $\tilde{\alpha}, \beta$ as in~\eqref{eq:alpha-beta},~\eqref{eq:alpha_tilde}. Then
\begin{equation*}
    \Tcal \ge \sum_{j=k+1}^{d}\nu_j^\downarrow - \frac{C}{g}\sbr{\rho^2 \sum_{i=1}^{k}\sum_{j=k+1}^{d}(\lambda_i + \lambda_j) + \sum_{j=1}^{k}\rbr{\nu_j^\downarrow}^2 + \rho^4 n\, k(d-k)} - \tilde{\Rcal}(\rho),
\end{equation*}
where the higher-order remainder $\tilde{\Rcal}(\rho)$ is the right-hand side of Lemma~\ref{lem:remainder_alter}.
\end{theorem}

\begin{corollary}\label{cor:smallcluster_alter}
There exist absolute constants $c_1, c_2, c_g > 0$ such that, if
\begin{equation*}
    \rho\sqrt{nd} \le c_1\,\sigma_k, \qquad \sigma_1 \le c_2\,\sigma_k, \qquad g > c_g \lambda_k,
\end{equation*}
then there is a constant $C > 0$ such that
\begin{equation*}
    \Tcal \ge \sum_{j=k+1}^{d}\nu_j^\downarrow - C\rbr{k(d-k)\,\rho^2 + \frac{k\,n^2 \rho^4}{\sigma_k^2} + \frac{d n^{3/2}\rho^3}{\sigma_k}}.
\end{equation*}
\end{corollary}

\begin{proof}

\emph{Verifying the hypotheses of Theorem~\ref{thm:smallcluster-main_alter}.} Using $\sigma_1 \le c_2\sigma_k$, $\rho\sqrt{nd} \le c_1\sigma_k$, and $n \ge d \ge 1$,
\begin{equation*}
    \sigma_1\rho\sqrt n \le \frac{c_2\sigma_k\cdot c_1\sigma_k}{\sqrt d} \le c_1 c_2\,\sigma_k^2,\qquad
    \rho^2 n = \frac{(\rho\sqrt{nd})^2}{d} \le c_1^2\,\sigma_k^2,
\end{equation*}
\begin{equation*}
    \sigma_1\rho\sqrt{nd} \le c_1 c_2\,\sigma_k^2,\qquad
    \rho^2 nd \le c_1^2\,\sigma_k^2.
\end{equation*}
Substituting into~\eqref{eq:alpha_tilde} and~\eqref{eq:alpha-beta}, and using $g \ge c_g\sigma_k^2$,
\begin{equation*}
    \tilde\alpha(\rho) \le \frac{2C(c_1 c_2 + c_1^2)}{c_g},\qquad
    \beta(\rho) \le \frac{2(2 c_1 c_2 + c_1^2)}{c_g}.
\end{equation*}
Choosing $c_1$ sufficiently small, we may ensure $\tilde\alpha(\rho), \beta(\rho) \le 1/2$ so that $\tfrac{1}{1-\tilde\alpha(\rho)},\ \tfrac{1}{1-\beta(\rho)} \le 2$. The hypotheses of Theorem~\ref{thm:smallcluster-main_alter} are thus satisfied.

\emph{Bounding the leading terms.} From $\sigma_1 \le c_2\sigma_k$, $\lambda_1 = \sigma_1^2 = O(\sigma_k^2)$ and hence $\lambda_i + \lambda_j = O(\sigma_k^2)$ for all $i \le k < j$. Therefore
\begin{equation*}
    \frac{\rho^2}{g}\sum_{i=1}^{k}\sum_{j=k+1}^{d}(\lambda_i+\lambda_j) \le C\,k(d-k)\,\rho^2.
\end{equation*}
Since each $\nu_j \le n\rho^2/4$,
\begin{equation*}
    \frac{1}{g}\sum_{j=1}^{k}(\nu_j^\downarrow)^2 \le C\frac{kn^2\rho^4}{\sigma_k^2}, \qquad \frac{\rho^4 n\,k(d-k)}{g} \le C\frac{\rho^4\,n\,k(d-k)}{\sigma_k^2} \le C\frac{kn^2\rho^4}{\sigma_k^2}.
\end{equation*}

\emph{Bounding the remainder.} Plugging~\eqref{eq:alpha_tilde} and~\eqref{eq:alpha-beta} into Lemma~\ref{lem:remainder_alter} together with $g = \Theta(\sigma_k^2)$ and $\lambda_1 = O(\sigma_k^2)$ gives
\begin{equation}\label{eq:R-cube-form-alter}
    \tilde{\Rcal}(\rho) \le \frac{Cd}{\sigma_k^4}\sbr{\rbr{\sigma_1\rho\sqrt n + \rho^2 n}^3 + 2 e^{-(n+d)}\rbr{2\sigma_1\rho\sqrt{nd}+\rho^2 nd}^3}.
\end{equation}
We claim the linear-in-$\rho$ part dominates the quadratic-in-$\rho$ part inside each cube. For the second cube, using $\sigma_1 \ge \sigma_k$ and $\rho\sqrt{nd} \le c_1\sigma_k$,
\begin{equation*}
    \frac{2\sigma_1\rho\sqrt{nd}}{\rho^2 nd} \;=\; \frac{2\sigma_1}{\rho\sqrt{nd}} \;\ge\; \frac{2\sigma_k}{c_1\sigma_k} \;=\; \frac{2}{c_1}.
\end{equation*}
For the first cube, write $\rho\sqrt n = \rho\sqrt{nd}/\sqrt d \le c_1\sigma_k/\sqrt d$. Using $\sigma_1 \ge \sigma_k$ and $d \ge 1$,
\begin{equation*}
    \frac{\sigma_1\rho\sqrt n}{\rho^2 n} \;=\; \frac{\sigma_1}{\rho\sqrt n} \;\ge\; \frac{\sigma_k}{c_1\sigma_k/\sqrt d} \;=\; \frac{\sqrt d}{c_1}.
\end{equation*}
(In contrast to Corollary~\ref{cor:smallcluster}, no separate lower bound on $d/\sqrt n$ is required: the alternative remainder uses $\sqrt n$ rather than $\sqrt d$ inside $\tilde\alpha$, so dominance follows directly from $\rho\sqrt{nd}\le c_1\sigma_k$.) Choosing $c_1 \le 1/2$, both ratios are $\ge 2$, so $(a+b)^3 \le 4 a^3$ for $a$ the linear part and $b$ the quadratic part. Thus $(\sigma_1\rho\sqrt n + \rho^2 n)^3 \le 4(\sigma_1\rho\sqrt n)^3$ and $(2\sigma_1\rho\sqrt{nd} + \rho^2 nd)^3 \le 4(2\sigma_1\rho\sqrt{nd})^3$. Substituting into~\eqref{eq:R-cube-form-alter} and using $\sigma_1 \le c_2\sigma_k$,
\begin{equation*}
    \tilde{\Rcal}(\rho) \le C\sbr{\frac{d\, n^{3/2}\rho^3}{\sigma_k} + 2 e^{-(n+d)}\,\frac{d^{5/2} n^{3/2}\rho^3}{\sigma_k}} \le C\,\frac{d\, n^{3/2}\rho^3}{\sigma_k},
\end{equation*}
where the last step absorbs the exponentially small bad-event contribution: since $d \le n$, $e^{-(n+d)} d^{3/2} \le e^{-n} n^{3/2} \to 0$ as $n \to \infty$.

Substituting these bounds into Theorem~\ref{thm:smallcluster-main_alter} proves the claim.
\end{proof}

\end{document}
